% easychair.tex,v 3.5 2017/03/15

\documentclass{easychair}

\usepackage{doc}

% use this if you have a long article and want to create an index
% \usepackage{makeidx}

% In order to save space or manage large tables or figures in a
% landcape-like text, you can use the rotating and pdflscape
% packages. Uncomment the desired from the below.
%
% \usepackage{rotating}
% \usepackage{pdflscape}

\makeatletter
\let\code\relax
\let\endcode\relax
\makeatother

\usepackage{agda}
\usepackage{newunicodechar}
\newunicodechar{⇒}{\ensuremath{\Rightarrow}}
\newunicodechar{◃}{\ensuremath{\triangleleft}}
\newunicodechar{▹}{\ensuremath{\triangleright}}
\newunicodechar{₁}{\ensuremath{_1}}
\newunicodechar{∈}{\ensuremath{\in}}
\newunicodechar{⊤}{\ensuremath{\top}}
\newunicodechar{≅}{\ensuremath{\cong}}

%\makeindex

%% Front Matter
%%
% Regular title as in the article class.
%
\title{Containers in Higher Kinds}

% Authors are joined by \and. Their affiliations are given by \inst, which indexes
% into the list defined using \institute
%
\author{
  Thorsten Altenkirch\inst{1} \and
  H{\aa}kon Robbestad Gylterud\inst{2} \and
  Zhili Tian\inst{1}
}

% Institutes for affiliations are also joined by \and,
\institute{
  $^{1}$School of Computer Science, University of Nottingham\\
  \email{\{psztxa,psxzt8\}@nottingham.ac.uk}\\
  $^{2}$Department of Informatics, University of Bergen\\
  \email{hakon.gylterud@uib.no}
}

%  \authorrunning{} has to be set for the shorter version of the authors' names;
% otherwise a warning will be rendered in the running heads. When processed by
% EasyChair, this command is mandatory: a document without \authorrunning
% will be rejected by EasyChair

\authorrunning{Altenkirch, Gylterud, Tian}

% \titlerunning{} has to be set to either the main title or its shorter
% version for the running heads. When processed by
% EasyChair, this command is mandatory: a document without \titlerunning
% will be rejected by EasyChair
\titlerunning{Containers in Higher Kinds}

\begin{document}

\maketitle

% \begin{abstract}
%   This is the abstract
% \end{abstract}

% The table of contents below is added for your convenience. Please do not use
% the table of contents if you are preparing your paper for publication in the
% EPiC Series or Kalpa Publications series

% \setcounter{tocdepth}{2}
% {\small
% \tableofcontents}

%\section{To mention}
%
%Processing in EasyChair - number of pages.
%
%Examples of how EasyChair processes papers. Caveats (replacement of EC
%class, errors).

\begin{code}[hide]
{-# OPTIONS --guardedness --allow-unsolved-metas #-}
module main where
\end{code}

Strictly positive types can be represented as containers \cite{containers},
that is \texttt{S : Set} and a family of positions \texttt{P : S → Set} giving rise to a
functor \texttt{⟦ S ◃ P ⟧ : Set ⇒ Set}. Every container has an initial algebra
(The W-type) and a terminal coalgebra (The M-type). However, there are
inductive/coinductive types which do not fit the scheme of W-types/M-types.
An example of such a type is the inductive type \texttt{Bush}, defined as follows:

\begin{code}
data Bush (A : Set) : Set where
  [] : Bush A
  _::_ : A → Bush (Bush A) → Bush A
\end{code}

Even though \texttt{Bush A} is not a W-type, the type constructor \texttt{Bush} itself arises
as the initial algebra of a functor on a functor category\cite{bird1998nested,
bird1999generalised}. This motivates the development of a notion of containers
at higher kinds (such as \texttt{(* ⇒ *) ⇒ (* ⇒ *)} or \texttt{(* ⇒ *) ⇒ *}) capable of modelling
strictly positive functors over containers.

\subsection*{Second-Order Containers}

We introduce a recursive syntax \texttt{2Cont} by extending the \texttt{Cont} with a family
\texttt{PF : S → Set} of positions of \texttt{F} and an inductive continuation \texttt{RF : (s : S) → PF s → 2Cont}.
These, second order containers are of kind \texttt{(* ⇒ *) ⇒ (* ⇒ *)}. Hence, each of them gives rise to an endofunctor \texttt{⟦ S ◃ PX ◃ PF ◃ RF ⟧ : Cont ⇒ Cont}.

The \texttt{Bush} example can be specified using the following higher-order signature:
\texttt{H F X = 1 + X × F (F X)}. There is a sum involved and therefore two shapes
(\texttt{S = Bool}). The left shape is trivial with a constant, while in the right
shape, there is one position of \texttt{X} (\texttt{PX true = ⊤}) and one position of
\texttt{F} (\texttt{PF true = ⊤}). Finally, we need to model \texttt{H' F X = F X} which proceeds
recursively (\texttt{RF true tt = H'}).

We define a second-order least-fixpoint operator \texttt{2W : 2Cont → Cont} by
induction-recursion, thereby recovering \texttt{Bush} as the initial algebra of
\texttt{H}. However, attempts to define a greatest fixpoint operator
\texttt{2M : 2Cont → Cont} by coinduction-induction shows positivity issue.

\subsection*{Higher-Order Containers}

To continue the story of constructing strictly positive functors of functor
categories, we define a notion of higher container \texttt{HCont : Ty → Set}, where
\texttt{Ty} are just the types of simply typed $\lambda$-calculus with one base type
\texttt{*} which represents \texttt{Set}. \texttt{HCont} is just a special case of
\texttt{Nf : Con → Ty → Set} where \texttt{Con} are the contexts of simply typed
$\lambda$-calculus \texttt{HCont A = Nf • A}. We also use \texttt{Var : Con → Ty → Set}
for the typed de Bruijn variables. For brevity, we do not present the full
syntax. We remark that normal forms |Nf| are defined mutually with neutral terms
and spines, following the standard presentation of hereditary substitution for
simply typed $\lambda$-calculus\cite{keller2010normalization}. With this setting,
it is easy to see that \texttt{Set ≅ HCont *}, \texttt{Cont ≅ HCont (* ⇒ *)} and
\texttt{2Cont ≅ HCont ((* ⇒ *) ⇒ * ⇒ *)}. 

We can interpret every higher container as a function on types, \texttt{⟦\_⟧ :
HCont A → ⟦ A ⟧T} where \texttt{⟦\_⟧T  : Ty → Set₁} is the obvious
interpretation of types (with \texttt{⟦ * ⟧T = Set}). We believe that this can
be extended to heredetary functors.

However, using this semantics
we can show that there is no third order fixed-point operators, \texttt{3W}.
Consider the following higher order container
\footnote{This is the corresponding $\lambda$-term, the actual encoding is not very readable.}:

\begin{code}[hide]
infixr 20 _⇒_
data Ty : Set where
  * : Ty
  _⇒_ : Ty → Ty → Ty

data HCont : Ty → Set where
\end{code}

\begin{code}
C : HCont (((* ⇒ *) ⇒ *) ⇒ (* ⇒ *) ⇒ *)
C = {!   !}

CC : ((Set → Set) → Set) → (Set → Set) → Set
CC F G = G (F G)
\end{code}

If an external fixpoint operator \texttt{3W : HCont (((* ⇒ *) ⇒ *) ⇒ (* ⇒ *) ⇒ *) → HCont ((* ⇒ *) ⇒ *)}
exists, applying it to \texttt{C} leads to an internal fixpoint operator \texttt{intW : HCont ((* ⇒ *) ⇒ *)}.
According to the semantics, \texttt{intW : (Set ⇒ Set) ⇒ Set} gives a least fixpoint of any functor, but we know such thing does not
exist.

We would like to incorporate initial algebras and terminal coalgebras
at any type either by adding explicit fixpoint operators or by
allowing infinite terms (i.e. using a coinductively defined syntax)
In either case we will have to give up the naive semantics sketched
above. 

\section*{Higher containers as normalized $\lambda$-terms}

Furthermore, we show higher containers are normalized $\lambda$-terms
closed under arbitrary products and coproducts (maybe also under least-fixpoint and
greatest-fixpoint). That is we define a syntax:

\begin{code}[hide]
infixl 5 _▹_
data Con : Set where
  •   : Con
  _▹_ : Con → Ty → Con

data Var : Con → Ty → Set where

variable
  Γ : Con
  A B : Ty
\end{code}

\begin{code}
data Tm : Con → Ty → Set₁ where
  var : Var Γ A → Tm Γ A
  lam : Tm (Γ ▹ A) B → Tm Γ (A ⇒ B)
  app : Tm Γ (A ⇒ B) → Tm Γ A → Tm Γ B
  Π : (I : Set) → (I → Tm Γ A) → Tm Γ A
  Σ : (I : Set) → (I → Tm Γ A) → Tm Γ A
\end{code}

together with a normalization function \texttt{nf : Tm Γ A → Nf Γ A}.

We also show \texttt{(Con, Tms, Ty, Tm)} is a model of simply-typed
$\lambda$-calculus (STLC), where \texttt{Tms} is a list of \texttt{Tm}. Similarly, we hope
to prove that the normalized model \texttt{(Con, Nfs, Ty, Nf)} is also a model of STLC, where
\texttt{Nfs} is a list of \texttt{Nf}.

\bibliographystyle{plain}
\bibliography{references}

%------------------------------------------------------------------------------
\end{document}